\section*{Annex F --- First Playtest Report}
\addcontentsline{toc}{section}{Annex F --- First Playtest Report}

The following is a first-person account of the first playtest session, held in March 2026 at the Jean Monet debate club in Bucharest. The session gathered six senior debaters (aged 18--20). Design changes arising from this playtest are documented separately in Annex B.

\subsubsection{Setting and Participants}

The playtest involved six senior debaters from the Jean Monet debate club in Bucharest, all aged between 18 and 20. These were experienced debaters, though none had encountered the Inside Out format before. I facilitated the session as the Evaluator and observed throughout.

\subsubsection{The Brief}

The brief took approximately 40 minutes. The format was significantly different enough from traditional debate that this length of explanation was necessary even for experienced players. Participants had many questions, particularly about two elements: (1) how the Subject role works and the roleplay dimension, which was described as ``the newest and most different part of the format'', and (2) whether counter-argumentation between Value Defenders is permitted --- a question I found somewhat surprising in hindsight.

\subsubsection{Preparation Time}

During prep time, the players struggled to figure out what to prepare. This was the first indication that the format demands a fundamentally different kind of preparation than traditional debate. Unlike preparing a policy or value case for traditional formats, players needed to figure out how to approach the Subject's dilemma through the lens of a specific value.

\subsubsection{Rituals}

The roll-in rituals felt awkward and cringe to participants. The rituals themselves were lengthy and formalized, which created discomfort rather than appropriate atmosphere-setting. More effective would be a brief introduction by the Evaluator followed by a small gesture or shared chant that everyone performs together.

Notably, the players did recognize and appreciate that the rituals acknowledged the effort and role of the judge --- in this case shared between the Evaluator and Subject.

The calling rituals (used to introduce each speaker) were received positively and did not require revision.

The roll-out rituals also felt overly lengthy and became cringe. However, one participant noted they appreciated the moment of collective gratitude that the rituals attempted to create. A simpler version involving handshakes and brief role-specific phrases would likely work better.

\subsubsection{Speech Flow and Embodiment}

An interesting pattern emerged across the speeches: players began using World Schools parliamentary plural and third-person references early on, but speech by speech, they naturally shifted toward second-person address and self-identification with the Subject's values. One player would say ``we believe'' at the start of their first speech, but by the second round was saying ``I believe in this value'' and referring to the Subject directly as ``you.''

This raised an important design question: should the Value Defenders fully embody the voice of the value (speaking as if they are an inner voice of the Subject themselves), or should they position themselves as advisors speaking about the Subject in the second person? The natural drift toward the former suggests players intuitively moved toward stronger roleplay embodiment, but the rules should be clarified on this point.

\subsubsection{First Round of Value Defenders}

Players immediately grasped that they needed to argue for the importance and relevance of their assigned value to the dilemma. They picked up on the fact that the solutions offered could be complex, realistic, and nuanced --- distinct from traditional debate solutions.

However, the distinction between arguing for a value and proposing a decision remained fuzzy for many. The fact that they must ``pick a solution and not necessarily a value'' was confusing initially; the rules require clarification and possibly example mini-speeches illustrating this distinction.

\subsubsection{Argumentation Style}

The format demands a form of argumentation that is fundamentally different from traditional debate. Rather than constructing large-scale causal chains and impact scenarios, Value Defenders must argue for their values on their own terms --- their importance, their weight in the decision, their interaction with other values. This was hard for players: it is uncommon in World Schools or BP debate, and finding ``more things to say'' about values themselves (as opposed to practical outcomes) required a different cognitive mode.

One significant observation: players felt the Subject character needed to be more fully fleshed out. The richer the characterization --- the beliefs they hold, decisions they have made in the past, their personality --- the easier it would be for Value Defenders to construct arguments tailored to that specific person's identity and narrative consistency. Players noted that this reflects how people actually construct their own identities: by looking at the coherence of their past actions and trying to remain narratively consistent.

\subsubsection{Participant Experience}

Despite the challenges, or perhaps because of them, participants reported the format to be surprisingly liberating. They reported feeling ``a lot less anxiety'' than in traditional debate, while also finding the work ``actually very complex and interesting strategically.'' They noted that architecting solutions was intellectually demanding in a different way: every detail of the context became potential argumentative material rather than merely setting for claims about large-scale consequences.

One highly competitive player --- the type that ``usually wins a lot at competitions'' --- reported finding the format less exhilarating. They felt that the cooperative nature took away from the ``thrill of collegial competitive jabbing'' and the ``adrenaline of winning,'' making the overall experience feel ``anticlimactic.'' This was a useful signal: the format appears to function as an alternative to traditional debate rather than a replacement, appealing to different motivations than ego-driven competitive success.

\subsubsection{Recognized Format Distinctiveness}

Players articulated several ways in which the format plays differently:

\begin{itemize}
    \item Explaining long, mechanistic causal chains is less necessary because there is no need to think through large-scale consequences (though such thinking could still occur if the dilemma involved someone with power over others).
    \item The power to formulate solutions makes it easier to stay topical and makes the whole debate less vague --- the debate becomes about real dilemmas rather than abstract policy positions.
    \item It is less about proving large-scale third-order impacts and more about navigating genuine tensions between values.
\end{itemize}

Players reported speaking more slowly but with similar theatricality to traditional debate, contrary to expectations. Though some initially felt they spoke more monotonously, other participants who reported understanding that the Subject could be emotionally moved as well noted they were quite theatrical and expressive. Some said that were they to debate again, they would try to be more emotionally immersed, consciously shifting away from trying to be very ``sweaty'' toward winning.

\subsubsection{Reported Learning and Growth}

When asked what they believed they had learned, participants identified several insights:

\begin{itemize}
    \item The format teaches you to be more relaxed and to enjoy playing.
    \item It teaches you to think about the context of decisions rather than abstractly applying rules.
    \item It makes you more human and more honest about human nature.
    \item It engages with the identity of decision-making: the recognition that ``every team says valuable stuff and it is important more rather to see what is important to us.''
    \item This last point led to a striking reflection: that sometimes decisions are not universal, they are personal, and arguing that someone is morally wrong when deciding or believing something is a bit absurd.
\end{itemize}

Some participants felt the format is more suitable for beginners than for highly competitive debaters. I agree with the latter assessment but also believe that with a more complex motion, experienced debaters could recognize how the format could be played at higher skill levels.

\subsubsection{Participant Feedback on Rules}

Many participants requested the addition of Points of Information (POIs) --- brief interjections during speeches. This became a high-priority feature request.

The participants also noted that the Subject role needs explicit scaffolding during preparation time. Rather than leaving this open, it would be helpful to have them prepare a first speech in which they explain the dilemma, why it is so hard, and how all values are relevant. Giving the dilemma ``real gravitas'' --- proving to everyone that the decision is truly important and truly hard --- would set the stage more effectively.

\subsubsection{Observational Notes and Design Recommendations}

Several design improvements emerged from observation:

\begin{itemize}
    \item \textbf{Ritual for calling people:} A more natural, shorter ritual is needed between speeches to call in the next speaker.
    \item \textbf{Subject preparation:} Give the Subject team something concrete to do during prep time: prepare a first speech that explains the dilemma and its difficulty, establishing that all values are genuinely relevant.
    \item \textbf{Motion design:} Include a longer list of things the Subject believes about themselves and decisions they have made in the past. This supports more specific, character-tailored argumentation by the Value Defenders.
    \item \textbf{Speech instructions:} Make explicit instructions clearer, particularly on the distinction between defending a value and proposing a solution, and on whether Value Defenders should embody an ``inner voice'' of the Subject or position themselves as advisors.
    \item \textbf{Interstitial speech:} Consider whether an interstitial speech between the first and second rounds of Value Defenders might help reset the debate or provide needed structure. (This was later addressed in subsequent iterations by introducing the Huddle.)
\end{itemize}

%%%%%%%%%%%%%%%%%%%%%%%%%%%%%%%%%%%%%%%%%

\newpage
\section*{Annex G --- Second Playtest Report}
\addcontentsline{toc}{section}{Annex G --- Second Playtest Report}

The following is a first-person account of the second playtest session, held in April 2026. The session gathered seven veteran debaters. Design changes arising from this playtest are documented separately in Annex H.

\subsubsection{Setting}

I held the second playtest together with seven veteran debaters, close friends. We held it at the apartment that part of us commonly rent and share with other nerds for the single purpose of playing TTRPGs there. Because of that, the main room is dominated by a hilariously large, roughly round table. We held the debate seated at it, which one player noted created a council-feel to the whole session --- as opposed to a regular debate in which the seating mimics parliamentary discourse and thus feels profoundly more impersonal. I was not the Evaluator this time; I played as the second speaker of the A Inner Voice team.

\subsubsection{The Brief}

The brief took about 40 minutes --- similar to how much it took with the high-schoolers. The fact that it required that amount of time even with people accustomed to parsing lengthy rulebooks speaks to how much the design differs from traditional formats. The division of the ``judge'' role into two separate individuals (the Subject and the Evaluator) was noted as interesting by several players. The role of the Subject was by far the hardest for players to understand, especially the purpose of the opening Subject speech. I had, at the point of writing the rules, found it genuinely difficult to pin down what the Subject should do in that first speech. I wrote vague guidelines --- and this showed --- but in the end I believe it was for the better: I let the player decide what would be most interesting and useful. In that I was fortunate: the role was drawn by a player with substantial game-design experience, a professional L\&D background, many years as a trainer, and deep experience as a roleplayer.

\subsubsection{Rituals}

The roll-in rituals were experienced as cringe, and were surely too lengthy. Players felt they should be replaced by a short introduction from the Evaluator followed by a brief shared gesture or chant. This aligns with how regular debates open: the judge makes an introduction, then everyone bangs the table vigorously and the debate begins. One player noted that the rituals do recognise the effort put in by the judge --- in this case shared between the Subject and the Evaluator.

The calling rituals were received positively.

The roll-out rituals were also too lengthy and became cringe, but one player noted they liked the way in which people express gratitude to everyone. The consensus was that this could be replaced by a moment of handshakes in which each person says a short phrase based on their role. Players also noted that calling a defender by name during the calling ritual felt unnecessary.

\subsubsection{S0 --- The Opening Subject Speech}

The player who drew the Subject role delivered an interesting, roleplay-heavy speech --- they even played a soundtrack. They fleshed out the motion with many evocative details that further emotionally contextualised the scenario, then tried to explain how they care about all three values and why each is relevant to the decision. This in part exposed some nuanced personal biases that bled from the player into the character. Several players --- including that player --- noted that this institutionalises the already existing practice of debaters ``playing for the judge'': attempting to exploit the judge's known biases. I found this observation striking. The speech, although I was sceptical of it initially, does hold a great deal of value: fleshing out the scenario, giving a floor for surfacing biases. What it needs is clearer language for what players are supposed to do.

\subsubsection{First Round of Inner Voices}

All players immediately adopted a roleplay stance. They spoke directly to the Subject, addressing them by name and in the second person, looking at them, and choosing speech patterns that mirrored those the Subject had displayed. They also took the liberty to add small details to the motion --- to worldbuild.

I find this very promising but also somewhat concerning. Giving players unofficial worldbuilding powers may lead inexperienced players to twist the debate in an unhelpful direction. However, adding colour to the broad facts of the scenario does lend to a more realistic simulation of how people actually argue, raises emotional engagement, and heightens immersion.

It was interesting that even very veteran debaters struggled, especially in this first round, to both appeal to emotion and to argue specifically for their assigned value. Most focused on constructing interesting and nuanced Solutions and arguing why these were practical and useful --- utilitarian and cerebral overall --- except for the last team, who appealed explicitly to emotion, which set the tone interestingly for the second half.

Players also began referring to each other as ``the other inner voices'' --- an organic renaming that feels more natural than ``Value Defenders''. I plan to adopt it.

\subsubsection{The Huddle}

Everyone found the Huddle very, very fun. The Subject player later said it was a genuine eureka moment when they realised they could ask particular Inner Voices to have very specific conversations with each other. It is worth noting that this player has meeting facilitation experience. During the Huddle, the Subject first asked each team to clarify things they were still confused about, then asked teams to expand on arguments they found particularly interesting, then asked pairs of Inner Voices to directly clash on issues they felt needed resolving.

At a certain point the Subject ran out of ideas for what to ask. I believe that with two players on the Subject team this would happen less frequently. Still, ten minutes felt too long --- I would cut it to eight.

\subsubsection{Second Round of Inner Voices}

The Huddle exacerbated something that had already started in the first round: a slow drift toward consensus on what decision should be taken. The second round was much less about the practicalities of which decision is best, and much more about arguing which is the dominant lens through which the decision should be viewed.

My intuition is that something similar might happen with other groups, not only with veteran debaters.

During my own speech I also pivoted toward consensus on the right Solution. I did this because I found it interesting and novel --- something regular debate never allows --- and I wanted to see how it felt, to demonstrate that the rules permit it, and to explore whether I could deliberately fish for the ``Brave Pivot'' commendation. I found that I felt little need to rebut the points made by others; arguing for my lens and shifting focus felt like a more natural and persuasive strategy than attacking others, and I allowed myself to be openly persuaded by fellow Inner Voices on some points.

A fellow player later noted that in this format to outright attack what others say ``feels wrong'' --- disrespectful and unnecessary --- and that the format requires tact in how one goes about diminishing the relevance of another's points.

\subsubsection{Second Subject Speech}

The Subject player structured their speech quite like a regular debate adjudication: they explained what they understood each Inner Voice team to have argued and how convincing each case was, then expressed the solution they had reached and what had factored most in their decision. They later said that during the first speech they were in more of a roleplay mindframe, but that as the debate progressed they transitioned more and more to feeling and thinking like a debate judge.

My intention had been for the second speech to remain somewhat more roleplay-ish --- the Subject speaking about how something \textit{felt} persuasive, being a bit less analytical. I am not sure whether adjusting the instructions is needed, or whether this was simply the style of this particular player.

\subsubsection{The Evaluator Speech}

The Evaluator gave insightful feedback on how each of us could have been more persuasive to a regular person. She then noted it was very difficult to find feedback to give, because the format does not specify clearly toward what educational objectives the feedback should be directed.

This raised a broader question I should think about more: traditional debate exists within a complex, organic constellation of clubs and competitions. Most people tell themselves they do debate to learn --- critical thinking, public speaking --- and trainers and judges tell themselves they are there to teach useful skills. In practice, my sense is that the circuit has a life of its own, and that people invest effort on either side mostly for the social reputation it yields.

How should this format be played? Who should facilitate it? Who is it for? What should it help participants achieve?

Another interesting suggestion was that a structure allowing the Evaluator to be more actively involved --- to facilitate and give real-time feedback, to encourage players --- would be useful. I tend to agree. Renaming the Evaluator ``the Facilitator'' might help signal this.

The Evaluator also said that awarding commendations was difficult, and that in a larger circuit it would become easier: she would not feel pressure to give everyone a commendation and would reserve them for truly outstanding moments. She added that once the commendations were memorised the system would feel much more organic than the numeric speech-point quantification that most traditional formats require of judges.

\subsubsection{On Bleed}

I decided to skip the final bleed-releasing ritual, expecting it would feel awkward. It was interesting, then, that even without it, one of the players spontaneously shared during a brief intermission that the debate had hit uncannily close to home: on the way to the session, a family member had asked them to testify for them in court, which they had found morally wrong.

\subsubsection{On Commendations}

There was substantial discussion about the commendation system. The more ``gamer'' of the players expressed that it scratches a ``make a build'' itch, but felt that commendations should carry some bonuses as well. They also raised the concern that players might start farming them --- choosing to play in unusual ways simply to collect a specific commendation and ``complete'' a build.

\subsubsection{Participant Feedback}

Most players expressed that the game was quite fun and that they believed it would be enjoyable even outside a casual friends setting.

One noteworthy pattern in the feedback: women and players who had been less successful as traditional debaters said they strongly preferred this format. Men who had been highly successful debaters expressed finding it less exhilarating. One participant put it directly: ``There are few things more ego-boosting than the feeling of total mind-domination that you get when you win a debate, and this cannot give you that.''

One point everyone agreed on was that the format is much less anxiety-inducing than regular debate, while still providing a structure for meaningful conversation --- unlike open real-life discussions, where people tend to fight over talking space.

One player suggested that the format might benefit from transitioning more wholeheartedly toward improv theatre --- becoming a roleplaying game without any performance-evaluation component. Another countered that there is value in the sense that you are always progressing, at least a little, with every debate.

%%%%%%%%%%%%%%%%%%%%%%%%%%%%%%%%%%%%%%%%%

\newpage
\section*{Annex H --- Inside Out Debate Format: Third Design Iteration}
\addcontentsline{toc}{section}{Annex H --- Inside Out Debate Format: Third Design Iteration}

The third design iteration is based on observations and participant feedback gathered during the second playtest (April 2026). The changes below supersede the corresponding elements in the second iteration (Annex B).

\subsubsection{Role Renaming}

\textbf{Value Defenders are renamed Inner Voices.} This name emerged organically during the second playtest, when players began referring to each other as ``the other inner voices''. It better reflects the roleplay-oriented nature of the format and the relationship between the Value Defender teams and the Subject.

\subsubsection{S0 --- Opening Subject Speech Instructions}

The instructions for the opening Subject speech are revised to foreground the roleplay dimension:

\begin{quote}
Hold a monologue as if you are the character, alone in a place of comfort, speaking to yourself about the dilemma. Express why this decision is truly difficult and how much getting it right matters to you. Explain why each value is relevant to this decision, and how you cannot simply disregard any of them. Do not lean toward any decision yet.
\end{quote}

\subsubsection{The Huddle}

The Huddle is shortened from 10 minutes to \textbf{8 minutes}.

\subsubsection{S2 --- Final Subject Speech Instructions}

The instructions for the final Subject speech are revised to include emotional closure:

\begin{quote}
Synthesise all arguments and deliver a decision over a specific solution --- a concrete course of action from those presented by the Inner Voices. Explain how everything in the debate contributed to this decision; leave nothing important out. Say why you feel this is the right decision, and how you now feel about making it.
\end{quote}

\subsubsection{Seating}

A seating instruction is added: try to arrange chairs so that the Inner Voices and the Subject face each other. Inner Voices speak from where they sit, standing or seated. The Subject does the same. The Evaluator sits somewhere outside the main circle --- present but out of sight, out of mind.

\subsubsection{Specific Times}

The following timing guidelines are provided:

\begin{itemize}
    \item \textbf{Speaking times:} 5 minutes per speech.
    \item \textbf{Prep time:} 30 minutes for beginners or high-school debaters; 15 minutes for experienced adults.
    \item \textbf{Evaluator oral critique:} 10 minutes.
    \item \textbf{Evaluator-led debrief:} 15 minutes.
\end{itemize}

\subsubsection{Rituals}

The following rituals replace those described in the first design iteration (Annex A).

\paragraph{Rolling In.}
\begin{enumerate}
    \item \textbf{Evaluator opens} --- The Evaluator says: \textit{``You are here to decide together.''}
    \item \textbf{All respond} --- Everyone answers: \textit{``Yes, we are.''} followed by shared clapping.
\end{enumerate}

\paragraph{Calling an Inner Voice.}
\begin{enumerate}
    \item The Evaluator calls: \textit{``Voice, speak for [Value].''}
    \item The speaker \textbf{must begin their speech} with the words \textit{``I will\ldots''}
\end{enumerate}

\paragraph{Rolling Out.}
\begin{enumerate}
    \item \textbf{Voices' thanks} --- Each Inner Voice says: \textit{``Thank you for listening to me.''}
    \item \textbf{Subject's thanks} --- The Subject says to the Inner Voices: \textit{``Thank you for speaking your truth.''} The Subject then says to the Evaluator: \textit{``Thank you for listening to me.''}
    \item \textbf{Bleed check} --- Everyone states their real name and how they feel right now. The Evaluator then leads a brief debrief.
\end{enumerate}

\subsubsection{Commendation System}

The commendation system is updated with a tier structure governing how commendations are distributed:

\begin{itemize}
    \item Every debater should receive at least one \textbf{T1 (Earned)} commendation per speech. At higher levels of play, T1 commendations carry less competitive weight but remain a baseline acknowledgement of participation.
    \item The Evaluator has a \textbf{limited pool of T2 (Strong) and T3 (Exceptional)} commendations to offer per debate. These are awarded in private.
    \item Some commendations may be \textbf{exclusive to specific Evaluators}, reflecting their expertise or perspective.
    \item Some rounds or competitions may designate a \textbf{featured commendation} --- a particular commendation given special significance in that context.
\end{itemize}

\subsubsection{Evaluator Instructions}

The Evaluator role is reframed toward active facilitation:

\begin{itemize}
    \item You are the \textbf{facilitator} of the entire event. Your primary job is for everyone to feel safe and to play to grow.
    \item Feel free to \textbf{intervene} when it feels natural. Encourage players when you think it is necessary. Gently redirect when someone may be making others uncomfortable. Invite players out of their comfort zone when you believe it is in their best interest. Be human; establish a relationship of trust and growth with the others.
    \item After the debate ends, \textbf{first facilitate a brief debrief} stemming from the bleed check moment, then transition to the oral critique.
    \item Prepare \textbf{potential debrief questions} whilst others are preparing for the debate. Debrief questions should \textit{not} focus (as the oral critique does) on skill growth, but rather on \textbf{meaningful emotional growth and acceptance of others' realities}. Explore how debaters would act in a similar scenario and what their own values are regarding the values explored in the debate. When in doubt, use these two anchor questions: \textit{``How important are these values for you?''} and \textit{``Were you in a similar scenario, what other values would influence you?''}
    \item Open the oral critique with a \textbf{narrative of the debate} from your outside perspective: what each team did and how it all unfolded.
    \item Engage directly with each team: ask what they \textbf{felt they did well} and what they found difficult. Share your view and add your own observations.
    \item Ask each team what they want to \textbf{do better next time}, and give concrete advice on how to get there.
    \item Award at least one \textbf{T1 (Earned)} commendation per speech to every debater. Use your limited pool of \textbf{T2 (Strong) and T3 (Exceptional)} commendations sparingly --- these are awarded in private.
\end{itemize}

\newpage
